%% Elsevier single-column submission format for Digital Chemical Engineering
\documentclass[a4paper,fleqn]{cas-sc}

\usepackage[authoryear,longnamesfirst]{natbib}
\usepackage{amsmath}
\usepackage{amssymb}
\usepackage{capt-of}
\usepackage{float}
\usepackage{needspace}
\usepackage{lineno}
\graphicspath{{./}{docs/DCHE-D-26-00020/}}

%%% Author macros
\def\tsc#1{\csdef{#1}{\textsc{\lowercase{#1}}\xspace}}
\tsc{WGM}
\tsc{QE}

\ExplSyntaxOn
\keys_set:nn { stm / mktitle } { nologo }
\cs_set:Npn \__first_footerline:
{
    \group_begin:
    \small
    \sffamily
    \ifnum\theblind>0\relax
    \else
        \__short_authors:
    \fi
    \group_end:
}
\ExplSyntaxOff

\begin{document}
\let\WriteBookmarks\relax
\def\floatpagepagefraction{1}
\def\textpagefraction{.001}

\shorttitle{An Interpretable Statistical Surrogate of Activated Sludge Systems}
\shortauthors{}

\title[mode=title]{An Interpretable Statistical Surrogate of Activated Sludge Systems that Preserves Mass Conservation and Component Non-Negativity}

\author{}

\begin{abstract}
Activated sludge models provide a mechanistic basis for wastewater-process analysis, but their coupled non-linear kinetics limit the practical use of transparent surrogates for rapid design, analysis, and optimization. Standard machine-learning surrogates can improve predictive efficiency, yet they generally do not guarantee two basic physical requirements of steady-state activated sludge systems: exact conservation of the stoichiometrically implied material inventories and componentwise non-negativity of the predicted ASM state. This article develops Invariant-Constrained Second-Order Regression (ICSOR), an interpretable surrogate framework that addresses both requirements within a single steady-state formulation. ICSOR uses a partitioned second-order feature map over operational variables and influent ASM components, a learned coupled component-space driver, and deterministic recursive coupled-QP estimation to retain a transparent coefficient structure. At deployment, the method applies a hierarchical correction sequence consisting of a raw coupled prediction, an affine invariant restoration, and a final linear program that enforces exact conservation and non-negativity simultaneously. Because the model is formulated directly in ASM component space, reported effluent composites are obtained only through an external composition map after deployment. The study is developed for a standalone continuous stirred tank reactor and is positioned for comparison against unconstrained surrogate baselines to assess whether physically admissible predictions can be obtained without sacrificing interpretability or competitive predictive performance.
\end{abstract}

\begin{graphicalabstract}
\mbox{}
\end{graphicalabstract}

\begin{highlights}
\item \mbox{}
\item \mbox{}
\item \mbox{}
\end{highlights}

\begin{keywords}
Activated sludge \sep Surrogate modeling \sep Interpretable machine learning \sep Mass conservation \sep Non-negativity constraints \sep Second-order regression
\end{keywords}

\maketitle

\section*{Nomenclature}
\begin{tabbing}
\hspace{2.5cm} \= \kill
$A$ \> Full-row-rank invariant matrix defining the null space of the stoichiometric matrix \\
$B$ \> Second-order driver coefficient matrix \\
$c_{aff}$ \> Invariant-consistent affine correction used before constrained deployment \\
$c_{in}$ \> Influent ASM component concentration vector ($c_{in} \in \mathbb{R}^{F}$) \\
$c_{out}$ \> True steady-state effluent ASM component concentration vector ($c_{out} \in \mathbb{R}^{F}$) \\
$c_{raw}$ \> Unconstrained coupled component prediction before affine presolve \\
$\hat{c}$ \> Final deployed ICSOR ASM component prediction \\
$\widehat{C}$ \> Fitted non-negative component prediction matrix used in training \\
$D$ \> Total dimension of the second-order feature map \\
$F$ \> Number of mechanistic ASM components \\
$\Gamma$ \> Learned ASM component coupling matrix \\
$I_{comp}$ \> Composition matrix mapping ASM components to measured composites ($I_{comp} \in \mathbb{R}^{K \times F}$) \\
$K$ \> Number of measured composite variables \\
$\lambda_B$ \> Ridge regularization weight on $B$ \\
$\lambda_\Gamma$ \> Ridge regularization weight on $\Gamma$ \\
$\lambda_{inv}$ \> Invariant-residual penalty weight in the training objective \\
$\lambda_{sys}$ \> Coupled-system consistency penalty weight in the training objective \\
$M_{op}$ \> Number of operational variables \\
$N_{rxn}$ \> Number of reactions in the mechanistic model \\
$R$ \> Coupled-system matrix, $R = I_F - \Gamma$ \\
$u$ \> Operational input vector ($u \in \mathbb{R}^{M_{op}}$) \\
$w_c$ \> Componentwise weight vector used in deployed correction \\
$y_{ext}$ \> Reported measured effluent composite vector ($y_{ext} \in \mathbb{R}^{K}$) \\
$\hat{y}_{ext}$ \> Reported composite prediction, $\hat{y}_{ext} = I_{comp}\hat{c}$ \\
$\nu$ \> Stoichiometric matrix ($\nu \in \mathbb{R}^{N_{rxn} \times F}$) \\
$\xi$ \> Net reaction progress vector \\
$\phi$ \> Partitioned second-order feature map \\
$\otimes$ \> Kronecker product \\
\end{tabbing}

\section{Introduction}

Water pollution presents a critical environmental and public health challenge worldwide, with particularly severe repercussions in developing nations (Lin et al., 2022; Madhav et al., 2020). Biological wastewater treatment systems, particularly those based on the activated sludge process, are the dominant approach applied to combat this issue (Duan et al., 2022). However, their modeling and analysis present formidable engineering and computational challenges (Xiao et al., 2025). The underlying biophysical mechanisms are governed by the Activated Sludge Models (ASM) (Henze et al., 2006), which rely on coupled, multiplicative non-linear kinetic expressions to describe microbial growth and substrate consumption. This inherent non-linearity limits the practical application of transparent analytical methods for process design, frequently forcing engineers to rely on iterative trial-and-error approaches or computationally expensive heuristic algorithms (Aboagye et al., 2021; Padron-Paez et al., 2020).

In response to the computational and mathematical complexity inherent in non-linear Activated Sludge Models, research has increasingly focused on developing surrogate models to approximate the behavior of biological wastewater treatment systems (Bhosekar \& Ierapetritou, 2018; Durkin et al., 2024; Tsochatzidi et al., 2025). Traditional approaches have explored linearization techniques (Li et al., 2020; Yang et al., 2014) and multivariable polynomial fitting (Ahn et al., 2014; Kim et al., 2009; Lim et al., 2011). While linear methods like Partial Least Squares (PLS) offer interpretable mappings, they inherently struggle to capture the severe non-linearities of microbial kinetics without extensive feature engineering. To overcome these limitations, the machine learning (ML) paradigm has emerged as a powerful route for approximating highly non-linear wastewater-treatment behavior (Croll et al., 2023; Nazlabadi et al., 2021). Recent wastewater evidence further shows that, when predictive accuracy alone is the criterion, modern ML models can outperform mechanistic ASM baselines for effluent nitrogen prediction (Wang et al., 2025). A wide array of data-driven algorithms---ranging from distance-based methods like k-Nearest Neighbors (k-NN) and PLS (Khatri et al., 2020; Wold et al., 2001), to kernel methods like Support Vector Regression (SVR) (Fang et al., 2011), tree-based ensembles such as Random Forest (RF), AdaBoost, XGBoost, LightGBM, and CatBoost (Sadri Moghaddam \& Mesghali, 2023; Wang et al., 2025; Yang et al., 2025), and Artificial Neural Networks (ANN) (Asadi et al., 2017; Bagherzadeh et al., 2021; Mao et al., 2024)---have demonstrated strong predictive capabilities for effluent quality. 

Despite their predictive accuracy, these standard ML algorithms remain structurally blind to a central physical requirement of activated sludge systems: material cannot be created or destroyed arbitrarily across the modeled control volume. A surrogate may fit observed effluent behavior well while still implying an impossible redistribution of the underlying ASM components, thereby violating the mass balances that govern COD, nitrogen, or phosphorus transfer. Related composition-preserving ML studies show that even accurate tree-based predictors can exhibit persistent unphysical conservation residuals unless an explicit corrective mechanism is imposed (Sturm \& Silva, 2025). At the same time, physically admissible concentration states must remain componentwise non-negative. These two requirements are inseparable in practice. A surrogate that conserves mass but predicts negative concentrations is still unusable, and a surrogate that preserves non-negativity while violating mass conservation remains physically misleading.

Recent scientific-ML literature shows that satisfying both requirements simultaneously is nontrivial. A substantial class of methods enforces conservation alone through conservative flux formulations, stoichiometric architectures, projection operators, Bayesian reconciliation, or completion strategies, but does not guarantee non-negativity of the corrected state (Beucler et al., 2021; Geiss et al., 2022; Hansen et al., 2024; Harder et al., 2022; Sturm \& Silva, 2025; Sturm \& Wexler, 2020, 2022). In that broader literature, completion methods reconstruct withheld variables to satisfy linear equality constraints (Beucler et al., 2021), Bayesian reconciliation treats conservation as a probabilistic update (Hansen et al., 2024), and recent atmospheric-composition studies use stoichiometric null-space constructions and invariant-restoring nudges to keep learned states on the conservation manifold (Sturm \& Silva, 2025; Sturm \& Wexler, 2020, 2022). Closely related soft-constraint approaches, such as physics-informed neural networks, can encourage physically plausible behavior but generally provide neither exact conservation nor strict positivity guarantees; more generally, soft-penalty conservation treatments need not drive residual violations to zero at deployment (Hansen et al., 2024; Raissi et al., 2019). Positivity-only output transformations can prevent negative states, but without an accompanying conservation mechanism they do not by themselves enforce global material balance. Conversely, conservative projection steps that ignore positivity can still generate infeasible negative states in practice, which is precisely why joint enforcement remains difficult (Kircher \& Votsmeier, 2025).

A smaller body of work has reported simultaneous enforcement of conservation and non-negativity. Kircher and Votsmeier (2025) develop positivity-preserving conservative corrections for mechanistic kinetics surrogates in atmospheric chemistry and heterogeneous catalysis, while Sturm et al. (2023) show in aerosol transport that mass-preserving scaling can maintain both total conservation and non-negativity for advected aerosol superspecies. Min et al. (2024) and Utkarsh et al. (2025) likewise propose hard-constrained learning frameworks that can impose equality and inequality constraints together in more general machine-learning settings. These results demonstrate that dual physical guarantees are achievable, but they arise in application settings whose state definitions, constraints, and deployment objectives differ substantially from those of activated sludge process surrogates.

This gap is consequential for wastewater engineering. Activated sludge models couple multiple conserved material inventories with biologically and chemically constrained component states, so violations of either mass conservation or non-negativity can distort interpretation, confound process optimization, and undermine confidence in surrogate-assisted design. A physically constrained data-driven model would therefore offer several practical benefits: it would preserve the bookkeeping needed for COD, nitrogen, and phosphorus analysis, prevent impossible negative component concentrations, reduce reliance on ad hoc clipping or reconciliation, and yield predictions that remain more defensible as operating conditions and influent loads vary. Yet, to our knowledge, no prior activated sludge surrogate has provided hard guarantees of both stoichiometric mass conservation and componentwise non-negativity within a single interpretable steady-state framework.

This study addresses that need through the development of Invariant-Constrained Second-Order Regression (ICSOR), a physics-informed surrogate framework for steady-state activated sludge systems. ICSOR is designed to answer a practical question: given a steady-state influent condition and operating point, what effluent state should be predicted if the surrogate is required to remain physically admissible by construction? Rather than treating physical consistency as a soft regularizer or a purely diagnostic afterthought, the framework is built so that deployed predictions satisfy both mass-conservation requirements and non-negativity requirements simultaneously.

To do so, ICSOR couples an interpretable second-order surrogate with a hard physical correction performed in ASM component space. The regression structure retains transparent operating-condition, influent, and interaction effects, while the deployment step removes states that would violate the admissible physical region. The result is a surrogate that remains analytically structured and data-driven, yet is constrained to return only physically meaningful steady-state predictions.

The objective of this article is to formulate the ICSOR framework as an analytically structured surrogate and to evaluate its efficacy on a standalone Continuous Stirred Tank Reactor (CSTR). By comparing ICSOR against established unconstrained baselines, this study assesses whether the proposed model can achieve competitive predictive accuracy while explicitly guaranteeing the parsimony, interpretability, and rigorous physical consistency required for reliable process analysis. The framework presented herein is strictly bounded to steady-state predictions; while it guarantees mass conservation and non-negativity, it does not claim to enforce full kinetic or thermodynamic feasibility, nor does it replace dynamic trajectory simulators. Instead, it occupies a necessary middle ground: a surrogate simple enough to be estimated transparently from data, yet mathematically constrained to respect the fundamental physical constraints of wastewater engineering.

\section{Theory and Calculation}

\subsection{Physical Scope, State Spaces, and Notation}

A fixed reactor block or process unit represented by steady-state samples is considered. The system boundary used to define influent and effluent states is identical to the boundary used to define stoichiometric conservation. Any external source or sink that crosses that boundary must therefore be represented in the adopted stoichiometric model; otherwise, it lies outside the present conservation claim.

ICSOR is formulated natively in ASM component space. Its scientific output is the effluent component vector itself, not a measured aggregate. Measured composites are obtained only afterward through an external composition matrix. This choice is deliberate because mass conservation and componentwise non-negativity are both defined in ASM component coordinates.

Let $M_{op}$ be the number of operational variables, $F$ the number of ASM components, $K$ the number of reported composite variables, and $N_{rxn}$ the number of reactions. Single-sample vectors are written as column vectors. The following state spaces and operators are defined:
\begin{itemize}
    \item $u \in \mathbb{R}^{M_{op}}$: Operational input vector (e.g., hydraulic retention time, aeration intensity).
    \item $c_{in} \in \mathbb{R}^{F}$: Influent ASM component concentration vector.
    \item $c_{out} \in \mathbb{R}^{F}$: True steady-state effluent ASM component concentration vector.
    \item $\hat{c} \in \mathbb{R}^{F}$: Final deployed ICSOR component prediction.
    \item $y_{ext} \in \mathbb{R}^{K}$: Reported measured composite vector.
    \item $I_{comp} \in \mathbb{R}^{K \times F}$: Composition matrix mapping ASM component concentrations to reported composites, such that $y_{ext} = I_{comp}c_{out}$.
    \item $\nu \in \mathbb{R}^{N_{rxn} \times F}$: Stoichiometric matrix.
\end{itemize}

The linear composition map $I_{comp}$ remains external to the model itself. In the common wastewater-reporting case where its rows are entrywise non-negative, non-negative component predictions imply non-negative reported composites automatically. Thus ICSOR places the hard non-negativity guarantee where it is physically meaningful, namely on the predicted ASM component state.

\subsection{Stoichiometric Structure and Conserved Quantities}

Let $\nu \in \mathbb{R}^{N_{rxn} \times F}$ be the stoichiometric matrix written in the adopted ASM component basis. For a steady-state sample, define the net reaction progress vector $\xi \in \mathbb{R}^{N_{rxn}}$, scaled into concentration-equivalent units, such that the net change across the control volume is:
\begin{equation}
\label{eq:stoichiometric-change}
 c_{out} - c_{in} = \nu^T \xi
\end{equation}
The vector $\xi$ is rarely observed directly. To eliminate it and isolate the conserved quantities, a full-row-rank matrix $A \in \mathbb{R}^{q \times F}$ is introduced whose rows span the left null space of $\nu$, so that $A\nu^T = 0$. Multiplying the stoichiometric change relation by $A$ yields:
\begin{equation}
\label{eq:invariant-conservation}
 A(c_{out} - c_{in}) = A \nu^T \xi = 0 \quad \implies \quad A c_{out} = A c_{in}
\end{equation}
This invariant-matrix construction follows the same mathematical idea used by Sturm and Wexler (2020) in mass-conserving photochemical surrogates, where a null-space matrix is used to preserve elemental inventories exactly. In the present setting, the conserved rows of $A$ represent wastewater-material inventories implied by the adopted ASM reaction network and system boundary rather than atmospheric atom counts.
Each row of $A$ represents one independent conserved combination of ASM components implied by the adopted reaction network and system boundary. The basis of $A$ is not unique, but once it is fixed, the invariant relation $Ac_{out} = Ac_{in}$ defines the exact conservation law that deployed ICSOR predictions must satisfy.

\subsection{Coupled Second-Order Driver and Prediction-Space Training}

Operational variables $u$ alter the process environment, whereas influent concentrations $c_{in}$ determine the material inventory presented to that environment. ICSOR keeps these roles explicit by using a partitioned second-order feature map:
\begin{equation}
\label{eq:feature-map}
\phi(u, c_{in}) = \begin{bmatrix}
1 \\
 u \\
 c_{in} \\
 u \otimes u \\
 c_{in} \otimes c_{in} \\
 u \otimes c_{in}
\end{bmatrix} \in \mathbb{R}^{D}
\end{equation}
where $\otimes$ denotes the Kronecker product under column-wise vectorization and $D = 1 + M_{op} + F + M_{op}^2 + F^2 + M_{op}F$.

The feature map drives a latent component-space forcing vector rather than the final deployed prediction directly:
\begin{equation}
\label{eq:latent-driver}
 r(u, c_{in}) = B \phi(u, c_{in})
\end{equation}
Using a latent driver rather than predicting the deployed state directly is also consistent with stoichiometry-driven flux-to-tendency architectures such as the conservation-law neural-network formulation of Sturm and Wexler (2022). The present article adopts that separation in a different way: ICSOR keeps the driver explicitly linear-quadratic and interpretable, then combines it with a learned coupling matrix and a separate constrained deployment step rather than embedding a fixed stoichiometric layer inside a deep network.

Using the same partitioned basis, the driver can be written blockwise as
\begin{equation}
\label{eq:blockwise-driver}
 r(u, c_{in}) = b + W_u u + W_{in} c_{in} + \Theta_{uu}(u \otimes u) + \Theta_{cc}(c_{in} \otimes c_{in}) + \Theta_{uc}(u \otimes c_{in}),
\end{equation}
with $b \in \mathbb{R}^{F}$, $W_u \in \mathbb{R}^{F \times M_{op}}$, $W_{in} \in \mathbb{R}^{F \times F}$, $\Theta_{uu} \in \mathbb{R}^{F \times M_{op}^2}$, $\Theta_{cc} \in \mathbb{R}^{F \times F^2}$, and $\Theta_{uc} \in \mathbb{R}^{F \times (M_{op}F)}$.

Cross-component dependence is introduced explicitly through a learned coupling matrix $\Gamma \in \mathbb{R}^{F \times F}$. In this study, $\Gamma$ is estimated in an admissible set that fixes its diagonal to zero, bounds its off-diagonal entries, and rejects choices that make $R = I_F - \Gamma$ poorly conditioned. The coupled relation is written as
\begin{equation}
\label{eq:coupled-relation}
 R c = r(u, c_{in}), \qquad R = I_F - \Gamma
\end{equation}

The sign of $B$ is not constrained. That choice is deliberate: increasing and decreasing effects are both allowed in the latent driver, while non-negativity is enforced later where it matters physically, on the predicted component state itself.

Over a dataset of $N$ steady-state samples, let $\Phi \in \mathbb{R}^{N \times D}$ collect the feature rows, let $C_{in} \in \mathbb{R}_{+}^{N \times F}$ collect influent component states, and let $C_{out} \in \mathbb{R}_{+}^{N \times F}$ collect effluent component targets. ICSOR introduces a fitted non-negative component prediction matrix $\widehat{C} \in \mathbb{R}_{+}^{N \times F}$ and estimates $(B, \Gamma, \widehat{C})$ through the coupled objective
\begin{equation}
\label{eq:training-objective}
\begin{aligned}
(\widehat{B}, \widehat{\Gamma}, \widehat{C}) = {}&
\arg\min_{B,\, \Gamma \in \mathcal{G},\, \widehat{C} \in \mathbb{R}_{+}^{N \times F}}
\Bigl\|
C_{out} - \widehat{C}
\Bigr\|_F^2 \\
&+ \lambda_{inv}
\Bigl\|
(\widehat{C} - C_{in})A^T
\Bigr\|_F^2 \\
&+ \lambda_{sys}
\Bigl\|
\widehat{C}(I_F - \Gamma)^T - \Phi B^T
\Bigr\|_F^2 \\
&+ \lambda_B \|B\|_F^2
+ \lambda_\Gamma \|\Gamma\|_F^2.
\end{aligned}
\end{equation}

Each term has a distinct role. The first term fits effluent ASM components directly. The second penalizes invariant mismatch in component space. The third forces consistency between the fitted states and the coupled driver relation. The ridge penalties stabilize the driver and coupling estimates.

\subsection{Deployment-Time Inference with Exact Conservation and Non-Negativity}

Training determines $\widehat{B}$ and $\widehat{\Gamma}$. Deployment then predicts a new component state by solving a second mathematical program. For a new sample $(u_*, c_{in,*})$, define
\begin{equation}
\label{eq:deployment-definitions}
\phi_* = \phi(u_*, c_{in,*}), \qquad d_* = \widehat{B}\phi_*, \qquad \widehat{R} = I_F - \widehat{\Gamma}.
\end{equation}

Whenever $\widehat{R}$ is nonsingular, the unconstrained coupled state is
\begin{equation}
\label{eq:raw-coupled-state}
 c_{raw} = \widehat{R}^{-1} d_*.
\end{equation}

A cheap affine presolve then restores exact invariants without yet enforcing non-negativity:
\begin{equation}
\label{eq:affine-presolve}
 c_{aff} = \arg\min_{c \in \mathbb{R}^{F}} \; \|c - c_{raw}\|_2^2 \quad \text{subject to} \quad A c = A c_{in,*}.
\end{equation}
This affine presolve plays the role of a conservation-restoring nudge: like the species-weighted correction proposed by Sturm and Silva (2025), it returns the nearest invariant-consistent state to the raw predictor. Here, however, the affine correction is treated as a cheap intermediate stage rather than the final deployed output, because exact invariants alone do not preclude negative ASM component concentrations.

If $c_{raw}$ already satisfies the invariants and remains non-negative, ICSOR can deploy it directly. If $c_{raw}$ is infeasible but $c_{aff}$ is non-negative, the affine presolve is already sufficient. Only when the affine state still contains negative components does ICSOR activate the full correction program:
\begin{equation}
\label{eq:deployment-lp}
\begin{aligned}
\hat{c} = {}& \arg\min_{c \in \mathbb{R}^{F},\, \rho \in \mathbb{R}^{F}} \; w_c^T \rho \\
\text{subject to} \quad {}& A c = A c_{in,*}, \\
& c \ge 0, \\
& -\rho \le c - c_{aff} \le \rho, \\
& \rho \ge 0
\end{aligned}
\end{equation}
with fixed componentwise correction weights $w_c \in \mathbb{R}_{+}^{F}$. Equation~\eqref{eq:deployment-lp} is the deployed ICSOR predictor. It enforces exact invariant preservation and componentwise non-negativity simultaneously while keeping the final state as close as possible to the coupled affine reference in a weighted $L_1$ sense.
That second stage reflects the same central lesson emphasized by Kircher and Votsmeier (2025): a conservative correction is only physically useful if positivity is preserved at the same time. Their scaled-space formulation is specifically designed to keep small predicted concentrations from crossing into negative values during correction. The present adoption targets steady-state activated sludge component predictions rather than atmospheric chemistry or heterogeneous catalysis, and it uses a weighted $L_1$ correction about the affine invariant reference rather than the exact scaling strategy proposed there.

Under the adopted non-negative component basis, $c_{in,*}$ is feasible whenever the influent state itself is physically admissible. The deployed problem is therefore not an afterthought but the mechanism by which ICSOR closes the following gap: every deployed prediction preserves the adopted conserved quantities and remains non-negative componentwise. The raw-check, affine-presolve, then LP-correction sequence is also practically important because it reserves the general optimizer for the subset of samples whose non-negativity constraints are genuinely active.

\subsection{External Reporting and Blockwise Interpretability}

If reported measured composites are needed after deployment, they are computed externally by
\begin{equation}
\label{eq:external-reporting}
\hat{y}_{ext} = I_{comp}\hat{c}.
\end{equation}
This collapse is not part of the regression model itself; it is a reporting operator applied after deployment. Keeping it external matters because different internal ASM component states can produce the same measured aggregates. The learned scientific object is therefore $\hat{c}$, not the reported collapse alone.

Interpretability is retained at two levels. First, the six blocks of $B$ have direct engineering meaning in the latent driver: $b$ is the baseline forcing, $W_u$ and $W_{in}$ encode first-order operating and influent effects, $\Theta_{uu}$ and $\Theta_{cc}$ encode within-subspace curvature, and $\Theta_{uc}$ encodes operation-loading interaction effects. Second, $\Gamma$ describes how that driver is redistributed across ASM components through the coupled system matrix $R = I_F - \Gamma$. The deployed prediction must therefore be read as the result of a driver, a coupling structure, and a constrained inference solve acting together.

These coefficients are motivated by plant operation rather than by algebraic convenience. The operating block $W_u$ corresponds to levers that engineers actually move, such as aeration intensity and hydraulic retention time, while $W_{in}$ captures first-order loading carry-through from the influent state. The quadratic blocks then retain the regime dependence and interaction structure that plant interpretation requires: $\Theta_{uu}$ allows operating actions such as aeration to have different marginal effects across regimes, $\Theta_{uc}$ captures the fact that the same intervention acts differently under different influent loads, and $\Theta_{cc}$ captures coupled loading states such as simultaneous COD--nitrogen or phosphorus-related effects. In that sense, the block structure of $B$ remains tied to practical diagnoses of nitrification loss, carbon limitation, and phosphorus breakthrough rather than to a purely algebraic expansion. More generally, recent hard-constrained learning frameworks likewise emphasize that nonlinear interaction structure must be retained if one wants expressive surrogates after equality and inequality constraints are imposed (Utkarsh et al., 2025).

The coupling coefficients in $\Gamma$ serve a different purpose. They do not describe how the operating point generates forcing; they describe how that forcing is redistributed across ASM component balances once created. Couplings among dissolved oxygen and nitrifier channels, among ammonium, nitrite, and nitrate, among $S_{PO4}$, $X_{PAO}$, $X_{PP}$, and $X_{PHA}$, or between soluble and particulate COD-bearing states correspond to recurring material linkages that operators already recognize in nitrification, denitrification, phosphorus storage, and hydrolysis-uptake behavior. The sign of a coupling coefficient indicates whether that linkage is reinforcing or compensating, and its magnitude indicates how strongly the redistribution must be represented to reproduce the observed state pattern. Persistent entries of $\Gamma$ are therefore interpretable as recurring process linkages rather than as purely statistical artifacts.

\subsection{Deterministic Recursive Coupled-QP Estimation and Modeling Boundaries}

This article adopts the recursive coupled-QP training route rather than a generic stochastic first-order optimizer such as Adam. That choice is deliberate. For fixed $(\Gamma, \widehat{C})$, the $B$ update is a ridge-regularized linear least-squares solve. For fixed $(B, \widehat{C})$, the $\Gamma$ update is a convex box-constrained quadratic program. For fixed $(B, \Gamma)$, the $\widehat{C}$ update is a convex non-negative quadratic program with the invariant penalty retained. The training loop therefore consists of deterministic Gauss-Seidel sweeps over exact block solutions.

One practical reason is the matrix $R = I_F - \Gamma$ itself. Deployment requires $R$ to remain invertible, and in practice it must remain well conditioned because both the raw coupled state and the subsequent correction depend on solving through $R$. A nearly singular $R$ would amplify numerical noise, distort the affine presolve, and make the deployed correction unstable. The recursive coupled-QP route can control this requirement explicitly: each $\Gamma$ update is box-constrained, its diagonal is fixed at zero, and the candidate is accepted only if $\kappa(I_F - \Gamma)$ stays below a prescribed conditioning threshold. If that threshold is violated, the candidate can be shrunk toward the origin and, if necessary, replaced by a safe fallback. Stochastic methods such as Adam, and indeed other generic first-order solvers, do not provide that admissibility guarantee by themselves. They may lower the objective while still visiting or returning nearly singular coupling matrices unless an additional conditioning projection is imposed externally.

This architecture is preferred here over generic stochastic first-order training for three reasons. First, each exact block update cannot increase the coupled objective, so the objective sequence within a restart is monotone non-increasing and converges in value to a stationary point. Second, the optimization variables remain attached to interpretable convex subproblems instead of being absorbed into one optimizer trajectory with learning-rate, momentum, clipping, or epoch-count sensitivity. Third, with fixed initialization, restart policy, and solver tolerances, the recursive coupled-QP route is reproducible in a way that stochastic training need not be.

The solver choices follow the same principle of structural fidelity. OSQP is used for the convex QP subproblems because the $\Gamma$ and $\widehat{C}$ blocks are sparse, box- and non-negativity-constrained quadratic programs. The deployment problem in Eq.~\eqref{eq:deployment-lp} is solved with HiGHS because, once $c_{aff}$ is fixed, the deployed correction is a linear program rather than a quadratic one. A second computational advantage comes from the hierarchical presolve described in Section~2.4: the raw coupled state is checked first, the affine invariant correction is checked second, and the LP is reserved only for states whose non-negativity constraints remain active after those cheaper steps. The manuscript therefore presents one definite ICSOR contract: direct component-space prediction, explicit learned coupling, invariant-aware training, exact constrained deployment, and deterministic recursive coupled-QP estimation.

The iterative estimator and the accelerated deployment logic can be summarized as follows.

\Needspace{18\baselineskip}
\noindent\begin{minipage}{\textwidth}
\begin{quote}
\noindent\textbf{Pseudocode 1. Deterministic recursive coupled-QP training}
\begin{tabbing}
\hspace{0.6cm}\=\hspace{0.6cm}\=\hspace{0.6cm}\=\kill
\textbf{Input:} $\Phi$, $C_{in}$, $C_{out}$, $A$, solver settings \\
\textbf{for} restart $s = 1, \dots, S$ \textbf{do} \\
\> initialize admissible $\Gamma^{(0,s)}$ and non-negative $\widehat{C}^{(0,s)}$ \\
\> compute $B^{(0,s)}$ from the corresponding ridge solve \\
\> \textbf{for} $t = 0, \dots, T_{\max}-1$ \textbf{do} \\
\>\> update $B^{(t+1,s)}$ by ridge least squares \\
\>\> update $\Gamma^{(t+1,s)}$ by constrained QP \\
\>\> enforce zero diagonal, box bounds, and $\kappa(I_F - \Gamma^{(t+1,s)}) \le \kappa_{\max}$ \\
\>\> update $\widehat{C}^{(t+1,s)}$ by non-negative invariant-aware QP \\
\>\> record the objective value and its running minimum \\
\>\> \textbf{if} the trailing objective slope is below tolerance \textbf{then stop the restart} \\
\> \textbf{end for} \\
\textbf{end for} \\
retain the feasible restart with the smallest recorded objective \\
\textbf{Output:} $\widehat{B}$, $\widehat{\Gamma}$, $\widehat{C}$
\end{tabbing}
\end{quote}

\end{minipage}

\par\medskip

\Needspace{10\baselineskip}
\noindent\begin{minipage}{\textwidth}
\begin{quote}
\noindent\textbf{Pseudocode 2. Hierarchical deployment with affine presolve}
\begin{tabbing}
\hspace{0.6cm}\=\hspace{0.6cm}\=\hspace{0.6cm}\=\kill
\textbf{Input:} $u_*$, $c_{in,*}$, $\widehat{B}$, $\widehat{\Gamma}$, $A$ \\
form $\phi_*$ and compute the raw coupled state $c_{raw} = (I_F - \widehat{\Gamma})^{-1}\widehat{B}\phi_*$ \\
\textbf{if} $c_{raw}$ is invariant-consistent and non-negative \textbf{then return} $c_{raw}$ \\
form the affine invariant correction $c_{aff}$ \\
\textbf{if} $c_{aff}$ is non-negative \textbf{then return} $c_{aff}$ \\
solve the deployment LP in Eq.~\eqref{eq:deployment-lp} with $c_{aff}$ as the reference state \\
\textbf{return} the projected state $\hat{c}$
\end{tabbing}
\end{quote}
\end{minipage}

\section{Methodology}

\subsection{ASM2d-TSN Process Model, State Basis, and Reported Composites}

The benchmark was generated from a standalone steady-state continuous stirred tank reactor based on Activated Sludge Model No.\ 2d with two-step nitrification (ASM2d-TSN). The ASM2d-TSN implementation follows Guerrero et al. (2011), which was adopted as the basis for the process model; all stoichiometric values and simulation rate coefficients used here were taken from that study. The adopted process model contains $N_{rxn}=28$ coupled biological and chemical processes spanning aerobic, anoxic, and anaerobic hydrolysis; aerobic heterotrophic growth on fermentable substrate and acetate; heterotrophic denitrification through nitrate and nitrite pathways; fermentation; storage, growth, and lysis pathways for phosphate-accumulating organisms (PAO); growth and lysis of ammonia-oxidizing bacteria (AOB) and nitrite-oxidizing bacteria (NOB); and precipitation-redissolution reactions for phosphorus-bearing metal solids. The state basis contains $F=20$ ASM component states, comprising the dissolved variables $S_O$, $S_F$, $S_A$, $S_{NH4}$, $S_{NO2}$, $S_{NO3}$, $S_{N2}$, $S_{PO4}$, $S_I$, and $S_{ALK}$ together with the particulate variables $X_I$, $X_S$, $X_H$, $X_{PAO}$, $X_{PP}$, $X_{PHA}$, $X_{AOB}$, $X_{NOB}$, $X_{MeP}$, and $X_{MeOH}$.

For every simulated steady state, the reported outputs were computed externally from the component state through the same composition matrix $I_{comp}$ used throughout the study. The reported output basis therefore contains $K=4$ composites: chemical oxygen demand (COD), total nitrogen (TN), total phosphorus (TP), and total suspended solids (TSS). The composite calculations are not unique to ICSOR; they are part of the benchmark-data construction used for every model family.

Table~\ref{tab:simulation_domain} introduces the full Latin hypercube design used to generate the benchmark. The table is organized to distinguish the two manipulated operating variables from the soluble and particulate influent-component ranges, because those groups play different roles in the study: the operating variables define the reactor environment, whereas the influent components define the loading state presented to that environment. Read row-wise, the table gives the admissible interval for each of the 22 sampled inputs and therefore specifies the envelope within which the steady-state simulations, surrogate fitting, and later benchmarking claims are valid.

\begin{table}[htbp]
\centering
\caption{Ranges of the 22 sampled inputs used in the Latin hypercube design.}
\label{tab:simulation_domain}
\scriptsize
\setlength{\tabcolsep}{4pt}
\begin{tabular}{lll|lll}
\hline
Variable & Range & Unit & Variable & Range & Unit \\
\hline
$\mathrm{HRT}$ & [6.0, 36.0] & h & $S_I$ & [10.0, 90.0] & g COD/m$^3$ \\
$\mathrm{Aeration}$ & [0.5, 2.5] & -- & $S_{ALK}$ & [80.0, 260.0] & mol/m$^3$ \\
$S_O$ & [0.0, 0.5] & g O$_2$/m$^3$ & $X_I$ & [20.0, 120.0] & g COD/m$^3$ \\
$S_F$ & [20.0, 180.0] & g COD/m$^3$ & $X_S$ & [60.0, 280.0] & g COD/m$^3$ \\
$S_A$ & [5.0, 80.0] & g COD/m$^3$ & $X_H$ & [15.0, 100.0] & g COD/m$^3$ \\
$S_{NH4}$ & [12.0, 55.0] & g N/m$^3$ & $X_{PAO}$ & [5.0, 60.0] & g COD/m$^3$ \\
$S_{NO2}$ & [0.0, 3.0] & g N/m$^3$ & $X_{PP}$ & [2.0, 20.0] & g P/m$^3$ \\
$S_{NO3}$ & [0.0, 8.0] & g N/m$^3$ & $X_{PHA}$ & [1.0, 30.0] & g COD/m$^3$ \\
$S_{N2}$ & [0.0, 2.0] & g N/m$^3$ & $X_{AOB}$ & [0.5, 8.0] & g COD/m$^3$ \\
$S_{PO4}$ & [2.0, 18.0] & g P/m$^3$ & $X_{NOB}$ & [0.5, 8.0] & g COD/m$^3$ \\
 &  &  & $X_{MeP}$ & [0.0, 12.0] & g TSS/m$^3$ \\
 &  &  & $X_{MeOH}$ & [0.0, 12.0] & g TSS/m$^3$ \\
\hline
\end{tabular}
\end{table}

Within Table~\ref{tab:simulation_domain}, the first two rows define the controllable reactor settings and the remaining rows define the influent composition envelope. The soluble-state bounds cover oxygen, carbon, nitrogen, phosphorus, inert, and alkalinity loadings, while the particulate-state bounds cover inert solids, biodegradable substrate, biomass pools, storage polymers, and precipitated solids. That separation matters later when interpreting ICSOR coefficients, because the model treats operating conditions and influent components as distinct blocks in both feature construction and physical correction.

\subsection{Simulation Protocol, Initialization, and Steady-State Acceptance}

The 22-dimensional sampling domain, consisting of the two operating variables and the 20-component influent state, was explored by Latin hypercube sampling with random seed 42. Sampling continued until 10,000 valid steady states had been accepted. At every sampled operating-influent point, the steady-state ASM balance equations were solved under box constraints by nonlinear least squares with lower and upper state bounds of $10^{-8}$ and 1500, variable, residual, and gradient tolerances of $10^{-9}$, and a maximum of 1200 function evaluations. Up to 25 candidate attempts were permitted per target sample so that infeasible operating-influent draws could be replaced without changing the prescribed sample count.

The nonlinear solve was supplied with two physically informed initializations. The primary initialization blended the current influent state with the previous accepted steady state, promoted biomass pools that are known to persist across nearby operating conditions, and estimated dissolved oxygen from the hydraulic-dilution and oxygen-transfer balance. A second multistart initialization then applied stronger substrate depletion and stronger biomass enrichment to probe alternative feasible basins. Both initializations were clipped to the admissible state bounds before optimization. If neither initialization yielded a solution whose maximum absolute residual was no greater than $10^{-9}$, the model was dynamically relaxed for 20 days with a stiff backward differentiation formula integrator and the terminal state of that dynamic relaxation was reused as a new steady-state guess. Only converged steady states that satisfied the residual-acceptance threshold were retained, so the final dataset contains only biologically viable operating points.

Table~\ref{tab:initial_conditions} summarizes these two initialization routes in the exact order they were applied. The table should be read as a decision aid for the nonlinear solver: the primary column defines the first physically informed guess used at each sampled operating point, and the multistart column defines the more aggressive fallback used to probe an alternative feasible basin before dynamic relaxation is invoked. Each row isolates one part of that initialization logic so that the steady-state acceptance protocol can be reproduced without ambiguity.

\begin{table}[H]
\centering
\caption{Initialization heuristics used for the steady-state nonlinear solve.}
\label{tab:initial_conditions}
\scriptsize
\resizebox{\textwidth}{!}{%
\begin{tabular}{lll}
\hline
Quantity & Primary initialization & Multistart initialization \\
\hline
Base state & $0.55\,c_{out}^{(prev)} + 0.45\,c_{in}$ when a previous accepted state exists; otherwise influent-based & Primary initialization reused as base state \\
$S_A$ and $S_F$ & 70\% of the base value & 35\% of the primary value \\
$X_H$ & $\max(\text{base}, 0.25\,X_S)$ & $\max(\text{primary}, 0.45\,X_S)$ \\
$X_{PAO}$ & $\max(\text{base}, 1.2\,X_{PP})$ & $\max(\text{primary}, 1.8\,X_{PP})$ \\
$X_{AOB}$ & $\max(\text{base}, 0.03\,S_{NH4})$ & $\max(\text{primary}, 0.08\,S_{NH4})$ \\
$X_{NOB}$ & $\max(\text{base}, 0.10\,S_{NO2})$ & $\max(\text{primary}, 0.25\,S_{NO2})$ \\
$S_O$ & $\dfrac{D\,S_O^{(prev)} + k_La\,S_O^{sat}}{D + k_La}$ with $D=24/\mathrm{HRT}$, $k_La = 2.0 + 18.0\,\mathrm{Aeration}$, and $S_O^{sat}=8.5$ & Inherited from the primary initialization \\
Bounds & Clipped to $[10^{-6},1500]$ & Clipped to $[10^{-6},1500]$ \\
\hline
\end{tabular}%
}
\end{table}

The structure of Table~\ref{tab:initial_conditions} reflects the intended solver behavior. The base-state row carries forward information from the previous accepted steady state when that information exists, which improves continuity across neighboring Latin hypercube samples. The substrate and biomass rows then bias the initialization toward biologically plausible regimes by partially depleting readily consumed substrates and preserving key biomass pools. The dissolved-oxygen row supplies a separate process-based estimate tied to hydraulic dilution and oxygen transfer, and the bounds row ensures that every initialization remains inside the admissible numerical domain before the nonlinear least-squares solve begins.

\subsection{Constraint Operators, Benchmark Dataset, and Common Training Targets}

The invariant operator $A$ used by ICSOR was built from the same stoichiometric model that generated the dataset. Specifically, if $\nu \in \mathbb{R}^{28 \times 20}$ denotes the ASM2d-TSN Petersen matrix, then $A$ was obtained as a basis for the left null space of $\nu$ so that $A\nu^T=0$ and hence $Ac_{out}=Ac_{in}$ for admissible steady states. The null-space basis was obtained by singular value decomposition of $\nu^T$, after which the retained basis vectors were transposed into row form, rounded for numerical stability, thresholded to remove negligible coefficients, and normalized by the first non-zero coefficient in each row.

A single benchmark dataset was constructed from the accepted steady states. Every model used the same exact 22 physical inputs, namely hydraulic retention time, aeration level, and the 20 influent ASM component concentrations. Every model was trained against the same 20-component effluent targets in ASM component space, and the same external composition map was then used to compute COD, TN, TP, and TSS from those predicted component states. The benchmark therefore holds fixed the sampled operating conditions, the training targets, and the composite-reporting rule across all model families. The methodological distinction lies only in the internal structure used to estimate the component state: ICSOR predicts the 20-component effluent state under invariant and non-negativity constraints using $c_{in}$ as the invariant reference in training and deployment, whereas the unconstrained regressors map the same 22-dimensional input directly to an unconstrained 20-component effluent prediction.

This common component-space training target determines the comparison protocol. Because the classical regressors were fitted without explicit knowledge of the conservation and non-negativity constraints, their coefficients were not optimized inside the restricted feasible search space that defines ICSOR. One could, in principle, append post-prediction constraint treatments such as analytic projection, null-space completion, conservative reconciliation, Bayesian correction, or positivity-preserving conservative adjustment (Beucler et al., 2021; Hansen et al., 2024; Kircher \& Votsmeier, 2025; Min et al., 2024; Sturm \& Silva, 2025; Sturm \& Wexler, 2020, 2022; Utkarsh et al., 2025). Such a deployment protocol was not used here. Applying an external correction only after an unconstrained component-space fit would alter the benchmark by giving the classical regressors a post hoc feasibility filter that is not inherent to their training problem, whereas ICSOR imposes the physical constraints from training through deployment and therefore searches a much more restricted model class from the outset.

\subsection{Benchmark Models and Optuna Hyperparameter Design}

The benchmark comprised ICSOR and nine unconstrained regressors: XGBoost, LightGBM, CatBoost, AdaBoost, Random Forest, Support Vector Regression, $k$-Nearest Neighbors, Partial Least Squares, and one fully connected multilayer perceptron (MLP). The MLP baseline contained three hidden layers of widths 128, 64, and 32, used logistic activation, and mapped the 22-dimensional input vector to the 20-component effluent state.

Figure~\ref{fig:mlp_architecture} summarizes the retained MLP architecture so that its input stage, hidden layers, output stage, and dense inter-layer connectivity are visually distinguished from the other benchmark models. The Optuna configuration used to tune every retained model is summarized in Table~\ref{tab:optuna_design}. The selected ICSOR settings are then listed separately in Table~\ref{tab:optuna_icsor} because that model carries the additional coupled and constraint-aware hyperparameters that do not appear in the unconstrained baselines. For the remaining baselines, Table~\ref{tab:optuna_classical_a} groups the selected settings for the tree-based regressors, Table~\ref{tab:optuna_classical_b} groups the kernel, instance-based, and linear regressors, and Table~\ref{tab:optuna_mlp} reports the final MLP-specific configuration. Taken together, these tables document both the common tuning contract and the final model-specific settings used in the benchmark and repeated dataset-size study.

\begin{center}
\includegraphics[width=\textwidth]{figure1_mlp_architecture.pdf}
\captionof{figure}{Architecture of the retained fully connected MLP baseline. The input and output stages are shown as labeled rectangular blocks to distinguish them from the hidden-layer neurons, which are shown as circular vertices. Representative vertices are displayed for the 22-input stage, the three hidden layers of widths 128, 64, and 32, and the 20-component effluent output stage, while the edge pattern indicates dense inter-layer connectivity between successive layers.}
\label{fig:mlp_architecture}
\end{center}

All retained models were assigned Optuna studies under a common tuning design. Hyperparameter selection minimized validation mean squared error in the shared 20-component effluent prediction space using a seeded tree-structured Parzen estimator sampler, a no-pruning policy, and 100 trials per model. The authoritative 8,000-row training pool was first reduced to a 4,000-row tuning subset, corresponding to 50\% of the training pool, and that subset was then partitioned with a 20\% validation fraction, yielding 3,200 tuning-training rows and 800 tuning-validation rows. No wall-clock timeout was imposed. Once the Optuna studies were completed, the final selected hyperparameters were fixed and reused for the one-shot benchmark and the repeated dataset-size analysis.

\begin{table}[H]
\centering
\caption{Common Optuna design used for all retained models.}
\label{tab:optuna_design}
\scriptsize
\begin{tabular}{ll}
\hline
Setting & Value \\
\hline
Objective & Validation mean squared error \\
Optimization direction & Minimize \\
Sampler & Tree-structured Parzen estimator \\
Pruning policy & None \\
Study seed & 42 \\
Trials per model & 100 \\
Training-pool fraction for tuning & 0.50 \\
Validation fraction within tuning subset & 0.20 \\
Tuning-training rows & 3,200 \\
Tuning-validation rows & 800 \\
Timeout & None \\
\hline
\end{tabular}
\end{table}

\begin{table}[H]
\centering
\caption{Optuna-selected ICSOR hyperparameters.}
\label{tab:optuna_icsor}
\scriptsize
\begin{tabular}{ll}
\hline
Parameter & Selected value \\
\hline
include\_bias\_term & true \\
$\lambda_{inv}$ & $8.68 \times 10^{-3}$ \\
$\lambda_{sys}$ & $1.27 \times 10^{-3}$ \\
$\lambda_B$ & $6.09 \times 10^{-5}$ \\
$\lambda_{\Gamma}$ & $6.55 \times 10^{-4}$ \\
$\gamma_{abs}$ & 0.368 \\
max\_outer\_iterations & 60 \\
n\_restarts & 1 \\
objective\_regression\_window & 12 \\
objective\_regression\_slope\_tolerance & $2.85 \times 10^{-7}$ \\
\hline
\end{tabular}
\end{table}

\begin{table}[H]
\centering
\caption{Optuna-selected hyperparameters for the tree-based regressors.}
\label{tab:optuna_classical_a}
\scriptsize
\begin{tabular}{>{\raggedright\arraybackslash}p{2.1cm}>{\raggedright\arraybackslash}p{9.0cm}}
\hline
Model & Selected settings \\
\hline
XGBoost & $n_{estimators}=309$, learning\_rate$=0.125$, max\_depth$=3$, min\_child\_weight$=6.84$, subsample$=0.915$, colsample\_bytree$=0.828$, reg\_alpha$=1.18 \times 10^{-3}$, reg\_lambda$=0.382$ \\
LightGBM & $n_{estimators}=302$, learning\_rate$=0.0809$, num\_leaves$=16$, max\_depth$=8$, min\_child\_samples$=14$, subsample$=0.845$, colsample\_bytree$=0.869$, reg\_alpha$=3.17 \times 10^{-3}$, reg\_lambda$=3.36 \times 10^{-4}$ \\
CatBoost & iterations$=236$, learning\_rate$=0.172$, depth$=4$, l2\_leaf\_reg$=0.0266$, bagging\_temperature$=0.0277$, random\_strength$=0.0657$ \\
AdaBoost & $n_{estimators}=236$, learning\_rate$=0.999$, loss$=$square \\
Random Forest & $n_{estimators}=390$, max\_depth$=14$, min\_samples\_split$=3$, min\_samples\_leaf$=1$, max\_features$=0.500$, bootstrap$=$false \\
\hline
\end{tabular}
\end{table}

\begin{table}[H]
\centering
\caption{Optuna-selected hyperparameters for the kernel, instance-based, and linear regressors.}
\label{tab:optuna_classical_b}
\scriptsize
\begin{tabular}{>{\raggedright\arraybackslash}p{2.1cm}>{\raggedright\arraybackslash}p{9.0cm}}
\hline
Model & Selected settings \\
\hline
SVR & $C=0.118$, epsilon$=0.0458$, kernel$=$poly, gamma$=$scale, degree$=4$, tol$=8.96 \times 10^{-4}$, coef0$=0.889$ \\
$k$-NN & n\_neighbors$=15$, weights$=$distance, algorithm$=$kd\_tree, leaf\_size$=40$, $p=2$, metric$=$manhattan \\
PLS & n\_components$=6$, max\_iter$=887$, tol$=9.91 \times 10^{-5}$ \\
\hline
\end{tabular}
\end{table}

\begin{table}[H]
\centering
\caption{Optuna-selected hyperparameters for the retained fully connected ANN.}
\label{tab:optuna_mlp}
\scriptsize
\begin{tabular}{ll}
\hline
Parameter & Selected value \\
\hline
hidden\_layer\_sizes & (128, 64, 32) \\
activation & logistic \\
solver & adam \\
alpha & $6.33 \times 10^{-4}$ \\
batch\_size & 64 \\
learning\_rate\_init & $1.08 \times 10^{-3}$ \\
max\_iter & 500 \\
tol & $1.63 \times 10^{-5}$ \\
shuffle & false \\
early\_stopping & false \\
\hline
\end{tabular}
\end{table}

\subsection{Evaluation Protocol}

A single row-wise 80--20 train-test split with random seed 42 was used for the one-shot benchmark, yielding 8,000 training rows and 2,000 test rows. The unconstrained regressors were standardized by z-score transformations fit on the training split only, separately for the 22-input feature matrix and the 20-output effluent-component matrix, and their predictions were inverse-transformed before external composite calculation and metric evaluation. ICSOR was kept in physical coordinates so that the invariant operator $A$, the coupled deployment rule, and the external reporting map $I_{comp}$ remained aligned with the formulation of Section~2. For all models, COD, TN, TP, and TSS were computed from the predicted component states through the same composition map, after which aggregate and per-target $R^2$, mean squared error, root mean squared error, mean absolute error, and mean absolute percentage error were reported over those four composites.

Only the final deployed component prediction of each model was used in the comparison. For ICSOR this is the constrained deployed component state, whereas for the classical regressors it is the direct unconstrained component prediction. Reported composite outputs were then obtained by applying the same external composition map to those final component states.

Sample-efficiency, computational scaling, and ranking stability were examined through a repeated dataset-size analysis. The 10,000-sample benchmark was subsampled without replacement at total sizes of 500, 1450, 2400, 3350, 4300, 5250, 6200, 7150, 8100, 9050, and 10,000 rows. Each total size was repeated 100 times with reproducible seeds derived from the base seed 42, and each repeated sample was re-split with the same 20\% test fraction, yielding 1,100 train-test re-estimations per model. The Optuna-selected settings listed above were held fixed throughout this sweep so that the resulting learning curves isolated sample-size sensitivity and split variability rather than retuning noise. For every fit within this sweep, wall-clock training time was recorded, and computational scaling with training size was summarized from the distribution of training seconds across the 100 repeats at each sample size. Predictive-performance ranking was based primarily on test RMSE computed from the final reported composite outputs and was supplemented by test MAE, MSE, MAPE, $R^2$, normalized area under the learning curve, dense average rank across target-metric-size combinations, the train-test RMSE gap at the largest training size, and per-target leaderboards. Writing $\bar{M}_{i,m}(n)$ for the repeat-averaged test metric of model $i$ at training size $n$ and metric $m$, the normalized area under the learning curve was computed by
\begin{equation}
\mathrm{nAUC}_{i,m} = \frac{1}{n_{\max} - n_{\min}} \int_{n_{\min}}^{n_{\max}} \bar{M}_{i,m}(n)\,dn,
\end{equation}
evaluated numerically by the trapezoidal rule over the sampled dataset sizes. For each metric $m$, target $k$, and size $n$, models were assigned a dense rank $r_{i,m,k}(n)$ from the repeat-averaged test metric, with lower-is-better ordering for RMSE, MAE, MAPE, and MSE and higher-is-better ordering for $R^2$; the metric-specific average rank and overall average rank were then
\begin{equation}
\bar{r}_{i,m} = \frac{1}{KS} \sum_{k=1}^{K} \sum_{s=1}^{S} r_{i,m,k}(n_s),
\qquad
\bar{r}_{i} = \frac{1}{|\mathcal{M}|} \sum_{m \in \mathcal{M}} \bar{r}_{i,m},
\end{equation}
where $\mathcal{M} = \{\mathrm{RMSE},\mathrm{MAE},R^2,\mathrm{MAPE},\mathrm{MSE}\}$. At the largest training size $n_{\max}$, the train-test RMSE gap was
\begin{equation}
\Delta_i^{\mathrm{RMSE}}(n_{\max}) = \overline{\mathrm{RMSE}}^{\mathrm{test}}_i(n_{\max}) - \overline{\mathrm{RMSE}}^{\mathrm{train}}_i(n_{\max}),
\end{equation}
so positive values indicate that the train split outperformed the test split. The per-target leaderboards were constructed by ranking models separately for each reported target $k$ at $n_{\max}$ using the repeat-averaged test RMSE,
\begin{equation}
r_{i,k}^{\mathrm{RMSE}}(n_{\max}) = \mathrm{DenseRank}_i\!\left(\overline{\mathrm{RMSE}}^{\mathrm{test}}_{i,k}(n_{\max})\right),
\end{equation}
and reporting the corresponding mean, standard deviation, and rank.

Physical-admissibility diagnostics were evaluated on the same final component-space predictions, because mass conservation and non-negativity are defined before external composition collapse. For a final component prediction $\tilde{c}$, mass-conservation violation was quantified by $v_{mc} = \|A(\tilde{c} - c_{in})\|_2$, and non-negativity violation was quantified by $v_{nn} = \|\min(\tilde{c}, 0)\|_1$, where the minimum is taken componentwise. A sample was counted as violating mass conservation or non-negativity when the corresponding quantity exceeded a numerical tolerance of $10^{-10}$. For each model, violation frequency was reported as the percentage of test samples that violated each condition.

\section{Results and Discussion Outline}

This section is an implementation scaffold for the final Results and Discussion section that will follow the methodology. The subsection order, paragraph roles, tables, and figures are arranged to match the evaluation protocol above while staying compatible with the flat-folder submission rules in the local submission guideline. Every figure file referenced below is stored directly in this submission folder with a manuscript-safe flat filename. The numerical values in the sample tables and figures are illustrative layout values only and should be replaced in place by the final benchmark exports.

\subsection{Overall Benchmark Performance on the Fixed Train-Test Split}

\noindent\textbf{Paragraph 1: Framing the benchmark.} Open the final prose with one short paragraph that restates the fixed 80--20 benchmark at a high level, points the reader to Table~\ref{tab:results_benchmark_sample}, and reminds the reader that every model is compared in the same externally reported composite space. This paragraph should identify the primary comparison metric used to order models in the narrative, which should remain aggregate test RMSE unless the final rerun shows a materially different discriminator.

\begin{table}[H]
\centering
\caption{Illustrative layout for the fixed-split benchmark summary. Replace the sample values with the final 80--20 test metrics for the retained model set before submission.}
\label{tab:results_benchmark_sample}
\scriptsize
\begin{tabular}{lcccc}
\hline
Model & RMSE & MAE & $R^2$ & MAPE \\
\hline
MLP & 4.887 & 2.305 & 0.9854 & 0.0156 \\
SVR & 5.772 & 3.111 & 0.9791 & 0.0223 \\
LightGBM & 6.500 & 3.681 & 0.9768 & 0.0246 \\
CatBoost & 6.514 & 3.746 & 0.9742 & 0.0269 \\
XGBoost & 6.824 & 3.955 & 0.9733 & 0.0270 \\
ICSOR & 6.942 & 3.929 & 0.9724 & 0.0266 \\
PLS & 10.348 & 6.271 & 0.9436 & 0.0420 \\
AdaBoost & 15.151 & 9.867 & 0.8815 & 0.0726 \\
Random Forest & 16.266 & 10.036 & 0.8806 & 0.0693 \\
$k$-NN & 18.166 & 11.913 & 0.7956 & 0.0941 \\
\hline
\end{tabular}
\end{table}

\noindent\textbf{Paragraph 2: Aggregate ranking and headline comparison.} Use the next paragraph to identify the best aggregate performer, place ICSOR relative to the strongest unconstrained baseline, and quantify the primary gap in the headline metric. Table~\ref{tab:results_benchmark_sample} should carry this discussion directly because it already contains the ranking metric together with the supporting secondary metrics, so a separate ranking figure would only restate the same evidence.

\noindent\textbf{Paragraph 3: Target-specific interpretation.} The third paragraph should explain whether the aggregate ranking hides target-specific reversals across COD, TN, TP, and TSS, using Table~\ref{tab:results_target_sample} as the anchor. This is the right place to state whether ICSOR is especially competitive on a physically constrained target, whether the unconstrained baselines dominate the carbon or solids targets, and whether any ranking reversals are large enough to change the take-home message from the aggregate table.

\begin{table}[htbp]
\centering
\caption{Illustrative layout for the per-target RMSE table at the largest training size. Replace the sample values with the final reported-composite RMSE values from the retained benchmark run.}
\label{tab:results_target_sample}
\scriptsize
\begin{tabular}{lcccc}
\hline
Model & COD & TN & TP & TSS \\
\hline
MLP & 7.647 & 1.047 & 0.473 & 6.434 \\
SVR & 8.638 & 1.594 & 0.453 & 7.747 \\
LightGBM & 10.597 & 1.685 & 0.470 & 8.234 \\
CatBoost & 10.191 & 1.798 & 0.563 & 8.604 \\
XGBoost & 11.106 & 1.742 & 0.484 & 8.701 \\
ICSOR & 10.993 & 1.646 & 0.396 & 8.922 \\
PLS & 16.947 & 1.996 & 0.490 & 13.360 \\
AdaBoost & 27.406 & 2.821 & 1.449 & 17.011 \\
Random Forest & 31.979 & 3.985 & 2.328 & 16.419 \\
$k$-NN & 32.922 & 6.527 & 3.410 & 20.201 \\
\hline
\end{tabular}
\end{table}

\noindent\textbf{Paragraph 4: Transition to the repeated-size analysis.} Close the fixed-split subsection with one short transition paragraph that explains why the full argument cannot stop at one split. This paragraph should signal that the next subsection evaluates sample efficiency, ranking stability, and overfitting behavior across repeated resampling rather than revisiting one-shot benchmark numbers.

\subsection{Sample Efficiency, Ranking Stability, and Generalization}

\noindent\textbf{Paragraph 1: Learning-curve overview.} Open this subsection by explaining the repeated dataset-size protocol in one sentence and immediately pointing the reader to Figure~\ref{fig:results_learning_curve_sample}. This retained figure has a unique purpose in the outline: it is the only compact way to show how model separation evolves with sample size. The final prose should describe how quickly the leading models improve with sample size, whether ICSOR is more competitive in the low-data regime, and whether the ranking becomes more or less stable as the training set grows.

\begin{figure}
\centering
\includegraphics[width=0.92\textwidth]{figure2_rmse_learning_curve.pdf}
\caption{Illustrative layout for the repeated dataset-size learning-curve figure. Replace this sample asset with the final RMSE learning-curve export from the retained benchmark sweep.}
\label{fig:results_learning_curve_sample}
\end{figure}

\noindent\textbf{Paragraph 2: Compact ranking summary.} Use the second paragraph to interpret Table~\ref{tab:results_scaling_sample} rather than opening separate narrative threads for every metric. This table has a distinct role from Figure~\ref{fig:results_learning_curve_sample}: it compresses the multi-metric ranking, the final-size point estimate, and the generalization summary into one exact reference surface. The final prose should explain how these quantities together support or challenge the fixed-split ranking.

\begin{table}[htbp]
\centering
\caption{Illustrative layout for the repeated-size summary table. Replace the sample values with the final retained-model exports for the learning-curve summary used in the manuscript.}
\label{tab:results_scaling_sample}
\scriptsize
\begin{tabular}{lcccc}
\hline
Model & Final RMSE & RMSE nAUC & Sample Avg. Rank & $\Delta$RMSE \\
\hline
MLP & 4.887 & 6.350 & 2.40 & 2.105 \\
SVR & 5.772 & 8.832 & 3.60 & 2.075 \\
LightGBM & 6.500 & 8.369 & 4.60 & 2.769 \\
CatBoost & 6.514 & 8.983 & 5.80 & 1.286 \\
XGBoost & 6.824 & 8.893 & 6.80 & 1.865 \\
ICSOR & 6.942 & 8.117 & 3.80 & 0.172 \\
PLS & 10.348 & 14.336 & 5.40 & -0.206 \\
AdaBoost & 15.151 & 21.462 & 8.60 & 0.101 \\
Random Forest & 16.266 & 18.715 & 9.80 & 13.017 \\
$k$-NN & 18.166 & 21.854 & 10.60 & 18.166 \\
\hline
\end{tabular}
\end{table}

\noindent\textbf{Paragraph 3: Overfitting and stability.} The next paragraph should turn from average performance to generalization behavior. This is where the final text should interpret the train-test RMSE gap, identify which models appear to win by raw capacity alone, and explain whether ICSOR's structure yields a practically important stability advantage even when it does not win the absolute point-estimate ranking.

\noindent\textbf{Paragraph 4: Runtime scaling and cost interpretation.} Close the subsection by pointing the reader to Figure~\ref{fig:results_runtime_sample}, which is now a required part of the outline rather than an optional supplement. This figure has a distinct role from Figure~\ref{fig:results_learning_curve_sample} and Table~\ref{tab:results_scaling_sample}: it shows how the computational cost of each model grows across the repeated dataset-size sweep, making it possible to discuss whether a model is merely accurate, or accurate at a practical training cost. The final prose should explain whether ICSOR's recursive constrained estimation carries a meaningful runtime penalty, whether that penalty scales smoothly with dataset size, and whether the cost is justified by the gains in feasibility or low-data stability.

\begin{figure}
\centering
\includegraphics[width=0.92\textwidth]{figure3_runtime_learning_curve.pdf}
\caption{Illustrative layout for the runtime learning-curve figure. Replace this sample asset with the final runtime export from the retained repeated dataset-size benchmark sweep. The y-axis is plotted on a logarithmic scale so training-time differences remain readable across the full model set.}
\label{fig:results_runtime_sample}
\end{figure}

\subsection{Physical Admissibility and Deployment Behavior}

\noindent\textbf{Paragraph 1: Hard-constraint outcome.} Anchor this subsection with Table~\ref{tab:results_physical_sample}. The final prose should state clearly whether ICSOR achieves exact mass-conservation and non-negativity compliance on the benchmark test predictions and how often the unconstrained baselines violate either condition.

\begin{table}[htbp]
\centering
\caption{Illustrative layout for the physical-admissibility summary. Replace the sample values with the final violation frequencies reported on the final component-space predictions.}
\label{tab:results_physical_sample}
\scriptsize
\begin{tabular}{lcc}
\hline
Model & Mass-Conservation Violations (\%) & Non-Negativity Violations (\%) \\
\hline
MLP & 78.0 & 0.0 \\
SVR & 81.0 & 0.2 \\
LightGBM & 85.0 & 0.1 \\
CatBoost & 86.0 & 0.1 \\
XGBoost & 87.0 & 0.2 \\
ICSOR & 0.0 & 0.0 \\
PLS & 90.0 & 0.8 \\
AdaBoost & 93.0 & 1.5 \\
Random Forest & 96.0 & 2.1 \\
$k$-NN & 98.0 & 3.0 \\
\hline
\end{tabular}
\end{table}

\noindent\textbf{Paragraph 2: Why the violations matter.} Use the next paragraph to explain why even low-frequency violations remain unacceptable in the activated-sludge setting. Table~\ref{tab:results_physical_sample} should remain the sole main-text evidence panel here because the retained quantities are only two frequencies and are easier to compare exactly in tabular form than in a redundant bar chart. The final prose should connect these diagnostics to material bookkeeping, process interpretation, and optimization reliability rather than treating them as abstract machine-learning penalties.

\noindent\textbf{Paragraph 3: Deployment-stage interpretation.} The third paragraph should explain how the hierarchical deployment rule behaved in practice. Once the final rerun is complete, add one sentence on how often the raw prediction was already feasible, how often the affine presolve was sufficient, and how often the LP stage was activated. If that stage-frequency summary becomes too long for the main paper, move the detailed counts to supplementary material and keep only the deployment takeaway here.

\noindent\textbf{Paragraph 4: Accuracy-versus-feasibility synthesis.} Close the subsection by tying the admissibility diagnostics back to the benchmark metrics. This is the core argumentative paragraph for the paper: it should state whether the feasibility guarantee appears to require only a moderate loss in predictive accuracy and whether that tradeoff is acceptable for steady-state wastewater surrogate deployment.

\subsection{Interpretable Structure of the Learned ICSOR Model}

\noindent\textbf{Paragraph 1: COD $\Theta_{cc}$ figure anchor.} Open the interpretability subsection by pointing the reader to Figure~\ref{fig:results_interpretability_sample}. In the main manuscript, this retained figure should now be limited to the COD $\Theta_{cc}$ block only, because that single panel isolates the cross-influent interaction structure most directly tied to carbon-loading interpretation without repeating the full multi-block atlas in the main paper. The final prose should explain that the complete COD, TN, TP, and TSS coefficient atlases are presented in supplementary material and that the main text uses the COD $\Theta_{cc}$ panel as the representative interpretability anchor.

\begin{figure}
\centering
\includegraphics[width=0.94\textwidth]{figure4_icsor_structure.pdf}
\caption{Illustrative layout for the COD-specific $\Theta_{cc}$ heatmap. Replace this sample asset with the final COD $\Theta_{cc}$ panel selected for the retained manuscript narrative, and move the full COD, TN, TP, and TSS coefficient atlases to supplementary material.}
\label{fig:results_interpretability_sample}
\end{figure}

\noindent\textbf{Paragraph 2: COD $\Theta_{cc}$ interpretation.} Use the second paragraph to interpret the COD-specific $\Theta_{cc}$ heatmap in engineering terms. The final prose should identify the strongest COD cross-loading interactions, explain whether the dominant retained entries cluster around soluble-carbon, particulate-COD, biomass, or phosphorus-bearing states, and connect those patterns to the expected carbon-limitation and storage regimes of the reactor.

\begin{table}[htbp]
\centering
\caption{Illustrative layout for the COD blockwise retained-feature summary used to support the ICSOR regularization and feature-selection discussion. Replace the sample counts with the final retained-coefficient totals reported for each COD coefficient block.}
\label{tab:results_icsor_regularization_sample}
\scriptsize
\setlength{\tabcolsep}{4pt}
\begin{tabular}{p{3.0cm}p{2.0cm}p{2.0cm}p{2.0cm}}
\hline
Coefficient block & \makecell[l]{Total\\coeffs.} & \makecell[l]{Retained\\coeffs.} & \makecell[l]{\%\\Retained} \\
\hline
$b$ & 1 & 1 & 100.0 \\
$W_u$ & 2 & 2 & 100.0 \\
$W_{in}$ & 20 & 9 & 45.0 \\
$\Theta_{uu}$ & 4 & 3 & 75.0 \\
$\Theta_{uc}$ & 40 & 11 & 27.5 \\
$\Theta_{cc}$ & 400 & 58 & 14.5 \\
Overall summary & 467 & 84 & 18.0 \\
\hline
\end{tabular}
\end{table}

\noindent\textbf{Paragraph 3: Regularization and feature-selection summary.} Anchor the next paragraph with Table~\ref{tab:results_icsor_regularization_sample}. The final prose should argue that ICSOR is regularized first by its model structure: bias, operational main effects, influent main effects, operational curvature, operation-influent interactions, influent curvature, and residual coupling are separated into distinct coefficient blocks with fixed physical meaning and sharply different dimensionality. That decomposition already narrows the admissible hypothesis class before any penalty is tuned, so the table should be used to show that feature selection emerges as a blockwise retention pattern rather than as an arbitrary search over all possible dense interactions. If the fitted COD model still uses ridge penalties, note that the penalty stabilizes coefficient magnitude, but emphasize that the more important interpretability claim is architectural: only a limited subset of entries remains materially active in the higher-dimensional blocks, especially $\Theta_{uc}$, $\Theta_{cc}$, and $\Gamma$, because the ICSOR structure itself suppresses unnecessary interactions and concentrates attention on the physically meaningful ones.

\noindent\textbf{Paragraph 4: Main-text versus supplementary boundary.} Close the section by stating the compression rule for the final manuscript explicitly: keep only the COD $\Theta_{cc}$ panel and the short regularization/screening table in the main text, and move the complete COD, TN, TP, and TSS coefficient atlases together with any additional blockwise response-surface diagnostics to supplementary material. This paragraph should explain that the main paper uses the COD interaction surface as the representative interpretability example, while the supplement preserves full coefficient-atlas completeness for all reported composites.

\section*{References}

\noindent Aboagye, E. A., Burnham, S. M., Dailey, J., Zia, R., Tran, C., Desai, M., \& Yenkie, K. M. (2021). Systematic design, optimization, and sustainability assessment for generation of efficient wastewater treatment networks. \textit{Water, 13}(9), 1326.

\noindent Ahn, J. Y., Chu, K. H., Yoo, S. S., Mang, J. S., Sung, B. W., \& Ko, K. B. (2014). Determination of optimal operating factors via modeling for livestock wastewater treatment: Comparison of simulated and experimental data. \textit{International Biodeterioration \& Biodegradation, 95}, 46--54.

\noindent Asadi, A., Verma, A., Yang, K., \& Mejabi, B. (2017). Wastewater treatment aeration process optimization: A data mining approach. \textit{Journal of Environmental Management, 203}, 630--639.

\noindent Bagherzadeh, F., Mehrani, M. J., Basirifard, M., \& Roostaei, J. (2021). Comparative study on total nitrogen prediction in wastewater treatment plant and effect of various feature selection methods on machine learning algorithms performance. \textit{Journal of Water Process Engineering, 41}, 102033.

\noindent Beucler, T., Pritchard, M., Rasp, S., Ott, J., Baldi, P., \& Gentine, P. (2021). Enforcing analytic constraints in neural networks emulating physical systems. \textit{Physical Review Letters, 126}(9), 098302.

\noindent Bhosekar, A., \& Ierapetritou, M. (2018). Advances in surrogate based modeling, feasibility analysis, and optimization: A review. \textit{Computers \& Chemical Engineering, 108}, 250--267.

\noindent Croll, H. C., Ikuma, K., Ong, S. K., \& Sarkar, S. (2023). Reinforcement learning applied to wastewater treatment process control optimization: Approaches, challenges, and path forward. \textit{Critical Reviews in Environmental Science and Technology, 53}(20), 1775--1794.

\noindent Duan, S., Zhang, Y., \& Zheng, S. (2022). Heterotrophic nitrifying bacteria in wastewater biological nitrogen removal systems: A review. \textit{Critical Reviews in Environmental Science and Technology, 52}(13), 2302--2338.

\noindent Durkin, A., Otte, L., \& Guo, M. (2024). Surrogate-based optimisation of process systems to recover resources from wastewater. \textit{Computers \& Chemical Engineering, 182}, 108584.

\noindent Fang, F., Ni, B., Li, W., Sheng, G., \& Yu, H. (2011). A simulation-based integrated approach to optimize the biological nutrient removal process in a full-scale wastewater treatment plant. \textit{Chemical Engineering Journal, 174}(2--3), 635--643.

\noindent Geiss, A., Silva, S. J., \& Hardin, J. C. (2022). Downscaling atmospheric chemistry simulations with physically consistent deep learning. \textit{Geoscientific Model Development, 15}(17), 6677--6694.

\noindent Guerrero, J., Guisasola, A., \& Baeza, J. A. (2011). The nature of the carbon source rules the competition between PAO and denitrifiers in systems for simultaneous biological nitrogen and phosphorus removal. \textit{Water Research, 45}(16), 4793--4802.

\noindent Harder, P., Watson-Parris, D., Stier, P., Strassel, D., Gauger, N. R., \& Keuper, J. (2022). Physics-informed learning of aerosol microphysics. \textit{Environmental Data Science, 1}, e20.

\noindent Hansen, D., Maddix, D. C., Alizadeh, S., Gupta, G., \& Mahoney, M. W. (2024). Learning physical models that can respect conservation laws. \textit{Physica D: Nonlinear Phenomena, 457}, 133952.

\noindent Henze, M., Gujer, W., Mino, T., \& van Loosedrecht, M. (2006). Activated Sludge Models ASM1, ASM2, ASM2d and ASM3. \textit{Water Intelligence Online, 5}(0).

\noindent Khatri, P., Gupta, K. K., \& Gupta, R. K. (2020). A review of partial least squares modeling (PLSM) for water quality analysis. \textit{Modeling Earth Systems and Environment, 6}, 1419--1426.

\noindent Kim, M., Rao, A. S., \& Yoo, C. (2009). Dual Optimization Strategy for N and P Removal in a Biological Wastewater Treatment Plant. \textit{Industrial and Engineering Chemistry Research, 48}(13), 6363--6371.

\noindent Kircher, J., \& Votsmeier, M. (2025). Machine learning surrogate models for mechanistic kinetics: Embedding atom balance and positivity. \textit{The Journal of Physical Chemistry Letters, 16}(19), 4715--4723.

\noindent Li, M., Hu, S., Xia, J., Wang, J., Song, X., \& Shen, H. (2020). Dissolved Oxygen Model Predictive Control for Activated Sludge Process Model Based on the Fuzzy C-means Cluster Algorithm. \textit{International Journal of Control, Automation and Systems, 18}(9), 2435--2444.

\noindent Lim, J., Lee, Y., Ataei, A., \& Yoo, C. (2011). Exploration of dual optimal conditions for COD and nitrogen removal in an MBR. \textit{Asia-Pacific Journal of Chemical Engineering, 6}(3), 433--440.

\noindent Lin, L., Yang, H., \& Xu, X. (2022). Effects of Water Pollution on Human Health and Disease Heterogeneity: A Review. \textit{Frontiers in Environmental Science, 10}, 880246.

\noindent Madhav, S., Ahamad, A., Singh, A. K., Kushawaha, J., Chauhan, J. S., Sharma, S., \& Singh, P. (2019). \textit{Water pollutants: sources and impact on the environment and human health}. In \textit{Sensors in Water Pollutants Monitoring: Role of Material}, 43--62.

\noindent Mao, Z., Li, X., Zhang, X., Li, D., Lu, J., Li, J., \& Zheng, F. (2024). Optimization of effluent quality and energy consumption of aeration process in wastewater treatment plants using artificial intelligence. \textit{Journal of Water Process Engineering, 63}, 105384.

\noindent Min, Y., Sonar, A., \& Azizan, N. (2024). Hard-constrained neural networks with universal approximation guarantees. \textit{arXiv preprint arXiv:2410.10807}.

\noindent Nazlabadi, E., Niaragh, E. K., \& Moghaddam, M. R. A. (2021). A systematic and critical review of two decades' application of response surface methodology in biological wastewater treatment processes. \textit{Desalination and Water Treatment, 228}, 92--120.

\noindent Padron-Paez, J. I., Almaraz, S. D. L., \& Rom\'an-Mart\'inez, A. (2020). Sustainable wastewater treatment plants design through multiobjective optimization. \textit{Computers \& Chemical Engineering, 140}, 106850.

\noindent Raissi, M., Perdikaris, P., \& Karniadakis, G. E. (2019). Physics-informed neural networks: A deep learning framework for solving forward and inverse problems involving nonlinear partial differential equations. \textit{Journal of Computational Physics, 378}, 686--707.

\noindent Sadri Moghaddam, S., \& Mesghali, H. (2023). A new hybrid ensemble approach for the prediction of effluent total nitrogen from a full-scale wastewater treatment plant using a combined trickling filter-activated sludge system. \textit{Environmental Science and Pollution Research, 30}(1), 1622--1639.

\noindent Sturm, P. M., Cappa, C. D., \& Wexler, A. S. (2023). Advecting superspecies: Efficiently modeling transport of organic aerosol with a mass-conserving dimensionality reduction method. \textit{Journal of Advances in Modeling Earth Systems, 15}(3), e2022MS003235.

\noindent Sturm, P. M., \& Silva, S. J. (2025). A nudge to the truth: Atom conservation as a hard constraint in models of atmospheric composition using a species-weighted correction. \textit{ACS ES\&T Air, 5}(1), 279--289.

\noindent Sturm, P. M., \& Wexler, A. S. (2020). A mass- and energy-conserving framework for using machine learning to speed computations: A photochemistry example. \textit{Geoscientific Model Development, 13}(9), 4435--4442.

\noindent Sturm, P. M., \& Wexler, A. S. (2022). Conservation laws in a neural network architecture: Enforcing the atom balance of a Julia-based photochemical model (v0.2.0). \textit{Geoscientific Model Development, 15}(9), 3417--3431.

\noindent Tsochatzidi, A., Cenci, F., Aroniada, M., \& Papageorgiou, L. G. (2025). A novel surrogate-based optimisation framework for pharmaceutical process systems. \textit{International Journal of Pharmaceutics, 682}, 125861.

\noindent Utkarsh, Maddix, D. C., Ma, R., Mahoney, M. W., \& Wang, Y. (2025). End-to-end probabilistic framework for learning with hard constraints. \textit{arXiv preprint}.

\noindent Wang, Y., Li, T., Bai, L., Yu, H., \& Qu, F. (2025). Comparison of interpretable machine learning models and mechanistic model for predicting effluent nitrogen in WWTP. \textit{Journal of Water Process Engineering, 77}, 108344.

\noindent Wold, S., Sj\"ostr\"om, M., \& Eriksson, L. (2001). PLS-regression: a basic tool of chemometrics. \textit{Chemometrics and Intelligent Laboratory Systems, 58}(2), 109--130.

\noindent Xiao, Y., Zhao, B., Zhang, X., \& Yue, X. (2025). A new paradigm in activated sludge management: an interpretable framework integrating mechanistic dynamics with a data-driven surrogate model. \textit{Environmental Research}, 123435.

\noindent Yang, T., Qiu, W., Ma, Y., Chadli, M., \& Zhang, L. (2014). Fuzzy model-based predictive control of dissolved oxygen in activated sludge processes. \textit{Neurocomputing, 136}, 88--95.

\noindent Yang, S. J., Kim, J., Eom, J., Kim, M., Lee, M., \& Lee, K. H. (2025). Machine learning based prediction of waste activated sludge generation for optimization of WWTP operational efficiency. \textit{Environmental Research}, 122990.

\end{document}